% --- Archon physics formalization source begin ---
% archon:physics
% archon:covers PhyXMiniProblems/problem_phyx_mini_0190.lean
% archon:source-report reports/phyx_mini/problem_phyx_mini_0190.source.json

\chapter{Physics problem phyx\_mini\_0190}
\label{ch:PhyXMiniProblems_problem_phyx_mini_0190}

\paragraph{Problem source.}
\#\# Physical scenario

A ship uses a sonar system (figure) to locate underwater objects.

\#\# Question

Find the wavelength of a \textbackslash{}( 262 \textbackslash{}, \textbackslash{}mathrm\{Hz\} \textbackslash{}) wave.

\#\# Answer choices

- A: A: 5.64m

- B: B: 6.45m

- C: C: 5.46m

- D: D: 6.54m

\#\# Figure caption (auxiliary; use the image as primary evidence)

The image depicts a scene involving wave interactions, possibly illustrating the Doppler effect. Key elements include:

1. **Objects**:
   - A ship at the bottom left with visible masts and sails, seemingly at rest on water or a surface.
   - An airplane or similar flying object at the top right with a pointed nose.

2. **Waves**:
   - Blue and black wave pattern radiating from the aircraft towards the ship.
   - The black waves are compressed nearer to the airplane and spread out toward the ship, representing wave propagation.

3. **Text and Symbols**:
   - A black arrow labeled "λ" (for wavelength) points between two black wavefronts.
   - A green arrow labeled "v" (for velocity) is oriented from the airplane, likely indicating motion direction. 

4. **Scene**:
   - The waves appear to originate from the airplane and travel towards the ship, demonstrating a relationship between the objects through wave interactions.

This setup suggests a demonstration or explanation of the Doppler effect, where the frequency of waves changes due to the relative motion between source and observer.

\#\# Dataset metadata

Wave Properties; Physical Model Grounding Reasoning, Multi-Formula Reasoning

\paragraph{Recorded answer/context.}
A: A: 5.64m

\paragraph{Figure/image path.}
/root/proposal\_for\_physic/science-mango/phyx\_mini\_run/phyx\_data/test\_image/190.png

\paragraph{Formalization target.}
create a compiling Lean file with sorry bodies at `PhyXMiniProblems/problem\_phyx\_mini\_0190.lean`. The Lean declarations must preserve the physical quantities, dimensions or dimensional roles, figure labels, governing-law hypotheses, and final relation expressed by this problem.
Use Mathlib/Physlib names found through LeanExplore where available. If a physics API is missing, introduce faithful local abstractions rather than scalar placeholder aliases.

\begin{theorem}[Physics formalization target]
\label{thm:physics:phyx_mini_0190:target}
\lean{PhyXMiniProblems.ProblemPhyXMini0190.problem_phyx_mini_0190}
\uses{def:physics:phyx-mini-0190:phyxminiproblems-problemphyxmini0190-lengthinmeters, def:physics:phyx-mini-0190:phyxminiproblems-problemphyxmini0190-sonarwavesetup, def:physics:phyx-mini-0190:phyxminiproblems-problemphyxmini0190-matchesproblemstatement, def:physics:phyx-mini-0190:phyxminiproblems-problemphyxmini0190-matchessuppliedfigure, def:physics:phyx-mini-0190:phyxminiproblems-problemphyxmini0190-hasphysicalsonarwaveparameters, def:physics:phyx-mini-0190:phyxminiproblems-problemphyxmini0190-usesstandardwatersoundspeed, def:physics:phyx-mini-0190:phyxminiproblems-problemphyxmini0190-satisfiesacousticwavekinematics, def:physics:phyx-mini-0190:phyxminiproblems-problemphyxmini0190-recordeddatasetanswer, def:physics:phyx-mini-0190:phyxminiproblems-problemphyxmini0190-isclosestdisplayedwavelengthchoice}
The assigned autoformalize agent should translate this physics problem into Lean declarations in the covered file, with theorem and lemma proof bodies written as `by sorry`.
\end{theorem}
\begin{proof}
This is an autoformalization task, not a proof task. Produce statements that can later be proved by the physics prover without weakening the physical model.
\end{proof}

% --- Archon declaration topology begin ---
\section{Declaration topology}
\begin{definition}[Frequency Quantity]
  \label{def:physics:phyx-mini-0190:phyxminiproblems-problemphyxmini0190-frequencyquantity}
  \lean{PhyXMiniProblems.ProblemPhyXMini0190.FrequencyQuantity}
  A nonnegative acoustic frequency, with inverse-time dimension
\end{definition}
\begin{proof}
  This is the declared model definition of Frequency Quantity.
\end{proof}
\begin{definition}[Length Quantity]
  \label{def:physics:phyx-mini-0190:phyxminiproblems-problemphyxmini0190-lengthquantity}
  \lean{PhyXMiniProblems.ProblemPhyXMini0190.LengthQuantity}
  A nonnegative wavelength, with length dimension
\end{definition}
\begin{proof}
  This is the declared model definition of Length Quantity.
\end{proof}
\begin{definition}[Speed Quantity]
  \label{def:physics:phyx-mini-0190:phyxminiproblems-problemphyxmini0190-speedquantity}
  \lean{PhyXMiniProblems.ProblemPhyXMini0190.SpeedQuantity}
  A nonnegative propagation speed, independent of the unit used to read it
\end{definition}
\begin{proof}
  This is the declared model definition of Speed Quantity.
\end{proof}
\begin{definition}[frequency Readout]
  \label{def:physics:phyx-mini-0190:phyxminiproblems-problemphyxmini0190-frequencyreadout}
  \lean{PhyXMiniProblems.ProblemPhyXMini0190.frequencyReadout}
  \uses{def:physics:phyx-mini-0190:phyxminiproblems-problemphyxmini0190-frequencyquantity}
  Read a physical frequency in inverse units of the selected time unit
\end{definition}
\begin{proof}
  This is the declared model definition of frequency Readout.
\end{proof}
\begin{definition}[frequency In Hertz]
  \label{def:physics:phyx-mini-0190:phyxminiproblems-problemphyxmini0190-frequencyinhertz}
  \lean{PhyXMiniProblems.ProblemPhyXMini0190.frequencyInHertz}
  \uses{def:physics:phyx-mini-0190:phyxminiproblems-problemphyxmini0190-frequencyquantity, def:physics:phyx-mini-0190:phyxminiproblems-problemphyxmini0190-frequencyreadout}
  Hertz readout, i.e. cycles per SI second
\end{definition}
\begin{proof}
  This is the declared model definition of frequency In Hertz.
\end{proof}
\begin{definition}[length Readout]
  \label{def:physics:phyx-mini-0190:phyxminiproblems-problemphyxmini0190-lengthreadout}
  \lean{PhyXMiniProblems.ProblemPhyXMini0190.lengthReadout}
  \uses{def:physics:phyx-mini-0190:phyxminiproblems-problemphyxmini0190-lengthquantity}
  Read a physical length in the selected length unit
\end{definition}
\begin{proof}
  This is the declared model definition of length Readout.
\end{proof}
\begin{definition}[length In Meters]
  \label{def:physics:phyx-mini-0190:phyxminiproblems-problemphyxmini0190-lengthinmeters}
  \lean{PhyXMiniProblems.ProblemPhyXMini0190.lengthInMeters}
  \uses{def:physics:phyx-mini-0190:phyxminiproblems-problemphyxmini0190-lengthquantity, def:physics:phyx-mini-0190:phyxminiproblems-problemphyxmini0190-lengthreadout}
  Meter readout of a wavelength
\end{definition}
\begin{proof}
  This is the declared model definition of length In Meters.
\end{proof}
\begin{definition}[speed Readout]
  \label{def:physics:phyx-mini-0190:phyxminiproblems-problemphyxmini0190-speedreadout}
  \lean{PhyXMiniProblems.ProblemPhyXMini0190.speedReadout}
  \uses{def:physics:phyx-mini-0190:phyxminiproblems-problemphyxmini0190-speedquantity}
  Read a speed in the selected length unit per selected time unit
\end{definition}
\begin{proof}
  This is the declared model definition of speed Readout.
\end{proof}
\begin{definition}[speed In Meters Per Second]
  \label{def:physics:phyx-mini-0190:phyxminiproblems-problemphyxmini0190-speedinmeterspersecond}
  \lean{PhyXMiniProblems.ProblemPhyXMini0190.speedInMetersPerSecond}
  \uses{def:physics:phyx-mini-0190:phyxminiproblems-problemphyxmini0190-speedquantity, def:physics:phyx-mini-0190:phyxminiproblems-problemphyxmini0190-speedreadout}
  Meter-per-second readout of the acoustic propagation speed
\end{definition}
\begin{proof}
  This is the declared model definition of speed In Meters Per Second.
\end{proof}
\begin{definition}[Locating System Kind]
  \label{def:physics:phyx-mini-0190:phyxminiproblems-problemphyxmini0190-locatingsystemkind}
  \lean{PhyXMiniProblems.ProblemPhyXMini0190.LocatingSystemKind}
  The physical locating system described in the prose
\end{definition}
\begin{proof}
  This is the declared model definition of Locating System Kind.
\end{proof}
\begin{definition}[Acoustic Medium]
  \label{def:physics:phyx-mini-0190:phyxminiproblems-problemphyxmini0190-acousticmedium}
  \lean{PhyXMiniProblems.ProblemPhyXMini0190.AcousticMedium}
  Propagation medium relevant to the requested acoustic wavelength
\end{definition}
\begin{proof}
  This is the declared model definition of Acoustic Medium.
\end{proof}
\begin{definition}[Acoustic Wave Kind]
  \label{def:physics:phyx-mini-0190:phyxminiproblems-problemphyxmini0190-acousticwavekind}
  \lean{PhyXMiniProblems.ProblemPhyXMini0190.AcousticWaveKind}
  Qualitative wave type used by the locating system
\end{definition}
\begin{proof}
  This is the declared model definition of Acoustic Wave Kind.
\end{proof}
\begin{definition}[Sonar Endpoint Role]
  \label{def:physics:phyx-mini-0190:phyxminiproblems-problemphyxmini0190-sonarendpointrole}
  \lean{PhyXMiniProblems.ProblemPhyXMini0190.SonarEndpointRole}
  Source and target roles in the prose's sonar scenario
\end{definition}
\begin{proof}
  This is the declared model definition of Sonar Endpoint Role.
\end{proof}
\begin{definition}[Figure Object Icon]
  \label{def:physics:phyx-mini-0190:phyxminiproblems-problemphyxmini0190-figureobjecticon}
  \lean{PhyXMiniProblems.ProblemPhyXMini0190.FigureObjectIcon}
  Object silhouettes actually visible in the supplied bitmap
\end{definition}
\begin{proof}
  This is the declared model definition of Figure Object Icon.
\end{proof}
\begin{definition}[Figure Quantity Label]
  \label{def:physics:phyx-mini-0190:phyxminiproblems-problemphyxmini0190-figurequantitylabel}
  \lean{PhyXMiniProblems.ProblemPhyXMini0190.FigureQuantityLabel}
  Mathematical symbols printed in the bitmap
\end{definition}
\begin{proof}
  This is the declared model definition of Figure Quantity Label.
\end{proof}
\begin{definition}[Figure Propagation Direction]
  \label{def:physics:phyx-mini-0190:phyxminiproblems-problemphyxmini0190-figurepropagationdirection}
  \lean{PhyXMiniProblems.ProblemPhyXMini0190.FigurePropagationDirection}
  Direction of the wave pattern and green arrow in the bitmap
\end{definition}
\begin{proof}
  This is the declared model definition of Figure Propagation Direction.
\end{proof}
\begin{definition}[Sonar Wave Figure]
  \label{def:physics:phyx-mini-0190:phyxminiproblems-problemphyxmini0190-sonarwavefigure}
  \lean{PhyXMiniProblems.ProblemPhyXMini0190.SonarWaveFigure}
  \uses{def:physics:phyx-mini-0190:phyxminiproblems-problemphyxmini0190-lengthquantity, def:physics:phyx-mini-0190:phyxminiproblems-problemphyxmini0190-speedquantity, def:physics:phyx-mini-0190:phyxminiproblems-problemphyxmini0190-figureobjecticon, def:physics:phyx-mini-0190:phyxminiproblems-problemphyxmini0190-figurequantitylabel, def:physics:phyx-mini-0190:phyxminiproblems-problemphyxmini0190-figurepropagationdirection}
  The figure-derived objects and labeled quantities. wavefrontSpacing and velocityArrowMagnitude are physical quantities, not scalar placeholders
\end{definition}
\begin{proof}
  This is the declared model definition of Sonar Wave Figure.
\end{proof}
\begin{definition}[Sonar Wave Setup]
  \label{def:physics:phyx-mini-0190:phyxminiproblems-problemphyxmini0190-sonarwavesetup}
  \lean{PhyXMiniProblems.ProblemPhyXMini0190.SonarWaveSetup}
  \uses{def:physics:phyx-mini-0190:phyxminiproblems-problemphyxmini0190-frequencyquantity, def:physics:phyx-mini-0190:phyxminiproblems-problemphyxmini0190-lengthquantity, def:physics:phyx-mini-0190:phyxminiproblems-problemphyxmini0190-speedquantity, def:physics:phyx-mini-0190:phyxminiproblems-problemphyxmini0190-locatingsystemkind, def:physics:phyx-mini-0190:phyxminiproblems-problemphyxmini0190-acousticmedium, def:physics:phyx-mini-0190:phyxminiproblems-problemphyxmini0190-acousticwavekind, def:physics:phyx-mini-0190:phyxminiproblems-problemphyxmini0190-sonarendpointrole, def:physics:phyx-mini-0190:phyxminiproblems-problemphyxmini0190-sonarwavefigure}
  The physical wave together with the prose roles and associated diagram. The wavelength is an independent unknown field and is not defined to be an answer-choice value
\end{definition}
\begin{proof}
  This is the declared model definition of Sonar Wave Setup.
\end{proof}
\begin{definition}[Matches Problem Statement]
  \label{def:physics:phyx-mini-0190:phyxminiproblems-problemphyxmini0190-matchesproblemstatement}
  \lean{PhyXMiniProblems.ProblemPhyXMini0190.MatchesProblemStatement}
  \uses{def:physics:phyx-mini-0190:phyxminiproblems-problemphyxmini0190-frequencyinhertz, def:physics:phyx-mini-0190:phyxminiproblems-problemphyxmini0190-sonarwavesetup}
  Facts stated in the prose. This predicate contains the given 262 Hz frequency but no numerical wavelength or answer label
\end{definition}
\begin{proof}
  This is the declared model definition of Matches Problem Statement.
\end{proof}
\begin{definition}[Matches Supplied Figure]
  \label{def:physics:phyx-mini-0190:phyxminiproblems-problemphyxmini0190-matchessuppliedfigure}
  \lean{PhyXMiniProblems.ProblemPhyXMini0190.MatchesSuppliedFigure}
  \uses{def:physics:phyx-mini-0190:phyxminiproblems-problemphyxmini0190-sonarwavesetup}
  Readout of the primary bitmap: an aircraft-like source icon faces a ship, the wavefront spacing is marked λ, and a green velocity arrow is marked v. The equalities identify the two pictured quantities with the modeled wave quantities but assign neither one a requested numerical answer
\end{definition}
\begin{proof}
  This is the declared model definition of Matches Supplied Figure.
\end{proof}
\begin{definition}[Has Physical Sonar Wave Parameters]
  \label{def:physics:phyx-mini-0190:phyxminiproblems-problemphyxmini0190-hasphysicalsonarwaveparameters}
  \lean{PhyXMiniProblems.ProblemPhyXMini0190.HasPhysicalSonarWaveParameters}
  \uses{def:physics:phyx-mini-0190:phyxminiproblems-problemphyxmini0190-frequencyinhertz, def:physics:phyx-mini-0190:phyxminiproblems-problemphyxmini0190-lengthinmeters, def:physics:phyx-mini-0190:phyxminiproblems-problemphyxmini0190-speedinmeterspersecond, def:physics:phyx-mini-0190:phyxminiproblems-problemphyxmini0190-sonarwavesetup}
  Positivity conditions selecting an ordinary propagating acoustic wave
\end{definition}
\begin{proof}
  This is the declared model definition of Has Physical Sonar Wave Parameters.
\end{proof}
\begin{definition}[Uses Standard Water Sound Speed]
  \label{def:physics:phyx-mini-0190:phyxminiproblems-problemphyxmini0190-usesstandardwatersoundspeed}
  \lean{PhyXMiniProblems.ProblemPhyXMini0190.UsesStandardWaterSoundSpeed}
  \uses{def:physics:phyx-mini-0190:phyxminiproblems-problemphyxmini0190-speedinmeterspersecond, def:physics:phyx-mini-0190:phyxminiproblems-problemphyxmini0190-sonarwavesetup}
  The standard textbook calibration used for the speed of sonar in water. The source does not print a speed, so this auxiliary environmental datum is kept separate from both the problem/figure observations and the wave law
\end{definition}
\begin{proof}
  This is the declared model definition of Uses Standard Water Sound Speed.
\end{proof}
\begin{definition}[Satisfies Acoustic Wave Kinematics]
  \label{def:physics:phyx-mini-0190:phyxminiproblems-problemphyxmini0190-satisfiesacousticwavekinematics}
  \lean{PhyXMiniProblems.ProblemPhyXMini0190.SatisfiesAcousticWaveKinematics}
  \uses{def:physics:phyx-mini-0190:phyxminiproblems-problemphyxmini0190-frequencyreadout, def:physics:phyx-mini-0190:phyxminiproblems-problemphyxmini0190-lengthreadout, def:physics:phyx-mini-0190:phyxminiproblems-problemphyxmini0190-speedreadout, def:physics:phyx-mini-0190:phyxminiproblems-problemphyxmini0190-sonarwavesetup}
  Wave kinematics v = f λ, expressed in every compatible choice of length and time units. It contains no problem-specific wavelength or answer value
\end{definition}
\begin{proof}
  This is the declared model definition of Satisfies Acoustic Wave Kinematics.
\end{proof}
\begin{lemma}[sonar Wavelength In Meters eq seven Hundred Forty div one Hundred Thirty One]
  \label{lem:physics:phyx-mini-0190:phyxminiproblems-problemphyxmini0190-sonarwavelengthinmeters-eq-sevenhundredforty-div-onehundredthirtyone}
  \lean{PhyXMiniProblems.ProblemPhyXMini0190.sonarWavelengthInMeters_eq_sevenHundredForty_div_oneHundredThirtyOne}
  \uses{def:physics:phyx-mini-0190:phyxminiproblems-problemphyxmini0190-lengthinmeters, def:physics:phyx-mini-0190:phyxminiproblems-problemphyxmini0190-sonarwavesetup, def:physics:phyx-mini-0190:phyxminiproblems-problemphyxmini0190-matchesproblemstatement, def:physics:phyx-mini-0190:phyxminiproblems-problemphyxmini0190-usesstandardwatersoundspeed, def:physics:phyx-mini-0190:phyxminiproblems-problemphyxmini0190-satisfiesacousticwavekinematics}
  At 262 Hz and 1480 m/s, the wave law gives the exact modeled wavelength 1480 / 262 = 740 / 131 meters. This is derived rather than included in any setup, figure, calibration, or law premise
\end{lemma}
\begin{proof}
  The conclusion follows from the governing relations and figure readouts named in the statement of sonar Wavelength In Meters eq seven Hundred Forty div one Hundred Thirty One.
\end{proof}
\begin{definition}[Answer Choice]
  \label{def:physics:phyx-mini-0190:phyxminiproblems-problemphyxmini0190-answerchoice}
  \lean{PhyXMiniProblems.ProblemPhyXMini0190.AnswerChoice}
  Labels of the four wavelength choices printed with the problem
\end{definition}
\begin{proof}
  This is the declared model definition of Answer Choice.
\end{proof}
\begin{definition}[displayed Answer Wavelength Meters]
  \label{def:physics:phyx-mini-0190:phyxminiproblems-problemphyxmini0190-displayedanswerwavelengthmeters}
  \lean{PhyXMiniProblems.ProblemPhyXMini0190.displayedAnswerWavelengthMeters}
  \uses{def:physics:phyx-mini-0190:phyxminiproblems-problemphyxmini0190-answerchoice}
  Wavelength printed beside an answer label, measured in meters
\end{definition}
\begin{proof}
  This is the declared model definition of displayed Answer Wavelength Meters.
\end{proof}
\begin{definition}[recorded Dataset Answer]
  \label{def:physics:phyx-mini-0190:phyxminiproblems-problemphyxmini0190-recordeddatasetanswer}
  \lean{PhyXMiniProblems.ProblemPhyXMini0190.recordedDatasetAnswer}
  \uses{def:physics:phyx-mini-0190:phyxminiproblems-problemphyxmini0190-answerchoice}
  The answer label recorded by the source dataset
\end{definition}
\begin{proof}
  This is the declared model definition of recorded Dataset Answer.
\end{proof}
\begin{definition}[Is Closest Displayed Wavelength Choice]
  \label{def:physics:phyx-mini-0190:phyxminiproblems-problemphyxmini0190-isclosestdisplayedwavelengthchoice}
  \lean{PhyXMiniProblems.ProblemPhyXMini0190.IsClosestDisplayedWavelengthChoice}
  \uses{def:physics:phyx-mini-0190:phyxminiproblems-problemphyxmini0190-lengthquantity, def:physics:phyx-mini-0190:phyxminiproblems-problemphyxmini0190-lengthinmeters, def:physics:phyx-mini-0190:phyxminiproblems-problemphyxmini0190-answerchoice, def:physics:phyx-mini-0190:phyxminiproblems-problemphyxmini0190-displayedanswerwavelengthmeters}
  A displayed option is a closest listed approximation to the modeled physical wavelength. This compares all four supplied values and does not define the unknown wavelength to equal any option
\end{definition}
\begin{proof}
  This is the declared model definition of Is Closest Displayed Wavelength Choice.
\end{proof}
% --- Archon declaration topology end ---
% --- Archon physics formalization source end ---
